\section{Exact period presentations and a corrected sign}\label{exact-period-presentations-and-a-corrected-sign}

For completeness we identify the source's Lemma 4.5 presentations,
rather than merely produce the same types and orders. Put

\[R_s=\{a^2,b^2,c^2,(ab)^6,x^{2^s},(ac)^2,
[a,x^2]^2,[a,x^4]\}.\]

For \(d=n-s\geq4\), \(d\ne5\), use the following extra relations:

\[\begin{array}{c|c|l}
d\bmod3&M&\text{period relations}\\\hline
0&2^{(d-3)/3}&x_1^M=x^{-2M},\quad y_1^M=x^{2M},\quad z_1^M=x^{-2M},\\
1&2^{(d-4)/3}&x_2^M=y_2^M=z_2^M=1,\\
2&2^{(d-5)/3}&x_3^M=x^{6M},\quad y_3^M=x^{-6M},\quad z_3^M=x^{6M}.
\end{array}\tag{14}\]

The first exponent is negative in (14). It is positive in the literal
v2 Lemma 4.5, physical page 12. All other entries agree with that
source. In the first and third branches \(M\) is even; in the middle
branch \(M\geq1\).

Here is a presentation proof valid for all parameters. Start with the
presentation \(R_s\) alone. The cyclic subgroup
\(C=\langle x^4\rangle\) is normal, has order at most \(2^{s-2}\), and
its quotient is (6). The affine group (6) retracts to \(D_8\), so the
presentation maps onto its compatible product with \(D_{2^{s+1}}\).
The kernel bound and the exact infinite-base version of Theorem AL.4
identify the two groups. In this product, each extra relation (14),
using (10), is precisely \((m v,1)=1\) for one of the three source
translation vectors \(v\). Normal closure kills exactly
\(m\Lambda_i\times1\). The lattice lies in the retraction kernel by
the parity condition already checked. Hence the quotient is exactly
the finite compatible product of Section 5. This proves its order,
marking and IP with no unproved extension consistency.

The sign correction is forced by an elementary identity, independent
of the affine calculation. For any involutions \(a,b,c\), the words (9)
satisfy

\[x_1^a=y_1^{-1}.\tag{15}\]

Indeed \(x_1=(b a b c b a)^2\), \(y_1=(b c b a b a)^2\), and conjugation by \(a\)
reverses the latter word. In the first branch the literal source imposes
both \(x_1^M=x^{2M}\) and \(y_1^M=x^{2M}\). Since \(M\) is even and
\([a,x^4]=1\), conjugating the first equation by \(a\) and comparing with
the inverse of the second yields

\[x^{4M}=1.\tag{16}\]

For \(n=10,s=4\), one has \(M=2\). Thus the literal presentation cannot
have the advertised rotation order sixteen. GAP computes order 1536
and rotation order eight for that presentation. With the corrected
first sign, its order is 3072, its rotation order is sixteen, and all
intersections hold. Three Lean declarations prove (15), the general
period-sign implication, and (16) directly from the literal conjugation
words. They do not rely on this numerical example.

This is a correction to an auxiliary construction in the bound edition,
not a failure of Theorem 3. The corrected family proves the required
existence statement for every parameter. It is not reported as a solved
open conjecture.
